\section{Preliminaries}
\label{sec:preliminaries}

\subsection{Pile sizes and logarithmic time}

Write $[N]=\{1,\ldots,N\}$ and $[k]_0=\{0,\ldots,k-1\}$. Throughout,
$k\ge2$, $N\ge2$, and all logarithms are natural. A cut array consists of
nonnegative integers $n_i^{(j)}$, $1\le j\le K$, $i\in[k]_0$, satisfying
$\sum_i n_i^{(j)}=N$. Put
\[
 D_k=\left\{p\in[0,1]^k:\sum_i p_i=1\right\},\qquad
 p^{(j)}=\frac1N(n_0^{(j)},\ldots,n_{k-1}^{(j)}),\qquad
 S_N=\frac K{\log N}.
\]
Column $j$ is placed at time $j/\log N$. Its contribution to an interval
with arbitrary endpoints is defined by linear interpolation. More
precisely, for a function $f$ on $D_k$, set
\begin{equation}\label{eq:empirical-integral}
 \int_r^s f(p_N(t))\,dt
 =\frac1{\log N}\sum_{j=1}^K
 \left[\bigl((s\log N-j+1)\vee0\bigr)\wedge1
       -\bigl((r\log N-j+1)\vee0\bigr)\wedge1\right]f(p^{(j)}).
\end{equation}
Equivalently, $p_N(t)=p^{(j)}$ on
$((j-1)/\log N,j/\log N)$; the integrand is extended by zero outside
$[0,S_N]$. The explicit formula specifies the convention at every
endpoint. In particular, on grid intervals this is exactly the normalized
sum of the corresponding column costs.

For $T\ge0$ we use $K=\lfloor T\log N\rfloor$. Then
\begin{equation}\label{eq:horizon-rounding}
 0\le T-S_N<\frac1{\log N}.
\end{equation}
The distinction between $T$ and $S_N$ will be retained in the estimates.

\subsection{Scalar costs and margins}
\label{subsec:scalar-costs}

For $p\in D_k$ and $\theta\ge1$, write
\[
 p_{\max}=\max_i p_i,\qquad a_p=\log\frac1{p_{\max}},\qquad
 \phi_p(\theta)=\sum_i p_i^\theta,\qquad
 \psi_p(\theta)=-\log\phi_p(\theta).
\]
Zero coordinates contribute zero to the power sum. The tilted
probabilities, cost, and entropy are
\begin{align}\label{eq:tilted-costs}
 q_i(p,\theta)&=\frac{p_i^\theta}{\phi_p(\theta)},\qquad
 \sum_i q_i(p,\theta)=1,\nonumber\\
 j_p(\theta)&=\sum_{i:p_i>0}q_i(p,\theta)\log\frac1{p_i},\qquad
 h_p(\theta)=-\sum_{i:q_i(p,\theta)>0}q_i(p,\theta)\log q_i(p,\theta).
\end{align}
Thus no logarithm of a zero coordinate is used in a nonzero summand.

Define the empirical costs by
\begin{align*}
 A^{(N)}(r,s)&=\int_r^s a_{p_N(t)}\,dt,&
 B^{(N)}_\theta(r,s)&=\int_r^s\psi_{p_N(t)}(\theta)\,dt,\\
 J^{(N)}_\theta(r,s)&=\int_r^s j_{p_N(t)}(\theta)\,dt,&
 H^{(N)}_\theta(r,s)&=\int_r^s h_{p_N(t)}(\theta)\,dt.
\end{align*}
The two empirical margins are
\begin{align}\label{eq:empirical-margins}
 \beta^{(N)}(s)&=1-A^{(N)}(0,s),\nonumber\\
 \Lambda_1^{(N)}&=\sup_{0\le s\le u\le S_N}\inf_{\theta\in[3,4]}
 \left\{\frac{\theta+3}{2}\beta^{(N)}(s)
       -B^{(N)}_\theta(s,u)-2B^{(N)}_2(u,S_N)\right\}.
\end{align}

A time profile is a family $(\mu_t)_{t\in\mathbb R}$ of probability
measures on $D_k$. No regularity in $t$ is implicit in this definition.
Whenever the scalar functions concerned are integrable, write
\[
 A(r,s)=\int_r^s\int_{D_k}a_p\,d\mu_t(p)\,dt,\qquad
 B_\theta(r,s)=\int_r^s\int_{D_k}\psi_p(\theta)\,d\mu_t(p)\,dt.
\]
Write $\beta(s)=1-A(0,s)$, so that $\Lambda_0(T)=\beta(T)$.
Recall the margins $\Lambda_0(s)=\beta(s)$ and
\begin{equation}\label{eq:macro-margin}
 \Lambda_1(T)=\sup_{0\le s\le u\le T}\inf_{\theta\in[3,4]}
 \left\{\frac{\theta+3}{2}\Lambda_0(s)
             -B_\theta(s,u)-2B_2(u,T)\right\}.
\end{equation}
The smaller expressions arising in the first-moment argument are
\begin{align}\label{eq:auxiliary-margins}
 \Lambda_2(T)&=\sup_{0\le s\le u\le T}\inf_{\theta\in[3,4]}
 \left\{\frac{\theta-1}{2}\Lambda_0(s)
             -B_\theta(s,u)-2B_2(u,T)\right\},\nonumber\\
 \Lambda_3(T)&=\sup_{0\le s\le T}\{\Lambda_0(s)-B_2(s,T)\}.
\end{align}
The supremum and infimum conventions in these definitions are literal;
their boundedness follows from the next lemma.

\begin{lemma}[One-column bounds]\label{lem:scalar-bounds}
For $p\in D_k$ and $2\le\theta\le4$,
\[
 0\le a_p\le\log k,\qquad
 0<\phi_p(\theta)\le1,\qquad
 0\le\psi_p(\theta)\le4\log k.
\]
Moreover, for every $\theta\ge1$,
\begin{equation}\label{eq:one-column-dominance}
 (\theta-1)a_p\le\psi_p(\theta).
\end{equation}
The function $a_p$ is continuous on $D_k$, and
$(p,\theta)\mapsto\psi_p(\theta)$ is continuous on $D_k\times[2,4]$.
\end{lemma}

\begin{proof}
We have $1/k\le p_{\max}\le1$, giving the bounds on $a_p$.
At least one coordinate is positive, so $\phi_p(\theta)>0$.
Since $p_i^\theta\le p_i$, the power sum is at most one. Choosing a
coordinate with $p_i\ge1/k$ gives
$\phi_p(\theta)\ge p_i^\theta\ge k^{-\theta}$ and hence the stated
upper bound on $\psi_p(\theta)$. Also,
\[
 \phi_p(\theta)=\sum_i p_i^{\theta-1}p_i
 \le p_{\max}^{\theta-1}\sum_i p_i=p_{\max}^{\theta-1}.
\]
Taking minus logarithms proves~\eqref{eq:one-column-dominance}; when
$\theta=1$ both sides are zero. Continuity of the maximum gives
continuity of $a_p$. Each summand $p_i^\theta$ is continuous on the
displayed domain, including at $p_i=0$, and the sum stays positive.
Composition with $-\log$ proves the last assertion.
\end{proof}

The same bounds hold after integration against $\mu_t$. In particular,
all spatial integrals above exist. Temporal integrability remains an
explicit hypothesis.

\begin{lemma}[Reduction and monotonicity of the margins]
\label{lem:margin-dominance}
Suppose $T\ge0$ and the functions
$t\mapsto\int a_p\,d\mu_t(p)$ and
$t\mapsto\int\psi_p(\theta)\,d\mu_t(p)$, $\theta\in[2,4]$, are
integrable on $[0,T]$. Then
\begin{equation}\label{eq:dominance-reduction}
 \Lambda_2(T)\le\Lambda_0(T),\qquad
 \Lambda_3(T)=\Lambda_0(T),\qquad
 \max_{0\le i\le3}\Lambda_i(T)=\max\{\Lambda_0(T),\Lambda_1(T)\}.
\end{equation}
Every $\Lambda_i$, $0\le i\le3$, is nonincreasing on $[0,T]$.
For $0\le S\le T$,
\begin{equation}\label{eq:lambda-horizon}
 0\le\Lambda_1(S)-\Lambda_1(T)\le8\log k\,(T-S).
\end{equation}
\end{lemma}

\begin{proof}
Integrating~\eqref{eq:one-column-dominance} gives
\[
 (\theta-1)A(r,s)\le B_\theta(r,s),\qquad 0\le r\le s\le T.
\]
Thus $\Lambda_0(s)-B_2(s,T)\le\Lambda_0(T)$; equality is obtained at
$s=T$, proving the assertion about $\Lambda_3$. In the infimum defining
$\Lambda_2$, use $\theta=3$. Its value is at most
\[
 \Lambda_0(s)-B_3(s,u)-2B_2(u,T)
 \le\Lambda_0(s)-2A(s,u)-2A(u,T)
 \le\Lambda_0(T).
\]
Taking the supremum proves~\eqref{eq:dominance-reduction}.

Nonnegativity of $a_p$ makes $\Lambda_0$ nonincreasing. To compare
$\Lambda_1(T)$ with $\Lambda_1(S)$, take an admissible pair $(s,u)$ at
time $T$ and replace it by $(s\wedge S,u\wedge S)$. All three intervals
of integration are shortened. Since their integrands are nonnegative
and $(\theta+3)/2\ge0$, the expression at time $T$ is no larger than
the expression at time $S$ for this clipped pair. Infima and then
suprema give $\Lambda_1(T)\le\Lambda_1(S)$. The same argument, with
$(\theta-1)/2$ in place of $(\theta+3)/2$, proves monotonicity of
$\Lambda_2$. The equality already proved handles $\Lambda_3$.

For the other direction in~\eqref{eq:lambda-horizon}, a pair admissible
at time $S$ is also admissible at time $T$, and the two corresponding
expressions differ by exactly $2B_2(S,T)$, independently of $\theta$.
Consequently
\[
 \Lambda_1(S)\le\Lambda_1(T)+2B_2(S,T)
 \le\Lambda_1(T)+8\log k\,(T-S).
\]
All uses of infima and suprema are legitimate: the expressions have
absolute value at most $4+20T\log k$ on the indicated domains, by
Lemma~\ref{lem:scalar-bounds}.
\end{proof}

\subsection{Local approximation of a time profile}
\label{subsec:profile-approximation}

A finite Borel partition $\mathcal P=\{D_1,\ldots,D_m\}$ of $D_k$
includes a representative $p^{[i]}\in D_i$ for every cell. We call it
$\chi$-regular if, for all $p,p'\in D_i$ and $\theta\in[2,4]$,
\begin{align}\label{eq:partition-regularity}
 |a_p-a_{p'}|&\le\chi,&
 |\psi_p(\theta)-\psi_{p'}(\theta)|&\le\chi,\nonumber\\
 p^{[i]}_l>2\chi&\ \Longrightarrow\
 p_l>\chi\ \text{ and }\left|\log\frac1{p_l}-\log\frac1{p^{[i]}_l}\right|\le\chi,
 \nonumber\\
 p^{[i]}_l\le2\chi&\ \Longrightarrow\ p_l\le3\chi.
\end{align}

\begin{lemma}[Regular partitions]\label{lem:regular-partition}
For every $0<\chi<1/k$, there is a finite $\chi$-regular Borel
partition of $D_k$.
\end{lemma}

\begin{proof}
Besides the two continuous functions in Lemma~\ref{lem:scalar-bounds},
consider the coordinate functions
$p\mapsto\log(1/(p_l\vee\chi))$. Compactness gives a common spatial
modulus of continuity, uniform in $\theta\in[2,4]$ and in the finitely
many coordinates. Partition $[0,1]^k$ into half-open coordinate boxes
of sufficiently small side length, and intersect with $D_k$, keeping
the nonempty intersections. Choose a representative in each cell.
The side length can be chosen below $\chi$, with oscillation of every
displayed continuous function at most $\chi$ on a box. If
$p^{[i]}_l>2\chi$, coordinate closeness gives $p_l>\chi$, so the
truncated logarithms are the actual logarithms. If
$p^{[i]}_l\le2\chi$, coordinate closeness gives $p_l\le3\chi$.
These are precisely~\eqref{eq:partition-regularity}.
\end{proof}

For a chosen partition, the scalar integrability condition on $[0,T]$
means that each of
\begin{equation}\label{eq:scalar-integrability}
 t\longmapsto\mu_t(D_i),\qquad
 t\longmapsto\int a_p\,d\mu_t(p),\qquad
 t\longmapsto\int\psi_p(\theta)\,d\mu_t(p)
 \quad(1\le i\le m,\ 2\le\theta\le4)
\end{equation}
is Lebesgue integrable on $[0,T]$.

For $\rho,\varphi>0$, an array with $K=\lfloor T\log N\rfloor$ is
$(\mathcal P,\rho,\varphi)$-almost like $(\mu_t)$ if, for every
$a\in\rho\mathbb Z$ with $0\le a<b=\min(a+\rho,T)$ and every cell,
\begin{equation}\label{eq:almost-like}
 \left|\frac{\#\{j:a\le j/\log N\le b,\ p^{(j)}\in D_i\}}{\log N}
       -\int_a^b\mu_t(D_i)\,dt\right|
 \le\varphi(b-a)+\frac1{\log N}.
\end{equation}
Closed endpoints in this counting condition are intentional. A column
is $g$-good if two distinct piles each contain at least $gN$ cards.
The local good-column condition on $[r,T]$, with parameters
$g,\gamma,\rho>0$, is
\begin{equation}\label{eq:local-good}
 \#\{j:a\le j/\log N\le a+\rho,\ p^{(j)}\text{ is }g\text{-good}\}
 \ge\gamma\rho\log N
 \quad(r\le a,\ a+\rho\le T).
\end{equation}
Here $a$ ranges over all real numbers in the indicated range, not only
the grid $\rho\mathbb Z$.

\begin{lemma}[Uniform scalar approximation]\label{lem:scalar-quadrature}
Let $T,\rho>0$, $\varphi,\omega,M\ge0$, and let
$\mathcal P=\{D_1,\ldots,D_m\}$ be a finite Borel partition with
representatives. Suppose $f:D_k\to[0,M]$ is continuous and has
oscillation at most $\omega$ on each cell. Assume integrability of
$t\mapsto\mu_t(D_i)$ and $t\mapsto\int f\,d\mu_t$ on $[0,T]$.
If~\eqref{eq:almost-like} holds, then, uniformly for
$0\le r\le s\le S_N$,
\begin{align}\label{eq:quadrature-bound}
 \left|\int_r^s f(p_N(t))\,dt-
              \int_r^s\int f\,d\mu_t\,dt\right|
 \le{}&2T\omega+Mm\varphi T+4M\rho\nonumber\\
 &+\frac{4Mm(T/\rho+2)+6M}{\log N}.
\end{align}
\end{lemma}

\begin{proof}
We first use the samples with $r<j/\log N\le s$. Passing from the
closed counting intervals in~\eqref{eq:almost-like} to these
half-open intervals removes at most one column. The error for the
count in each cell of a full time interval of length $\rho$ is
therefore at most $\varphi\rho+2/\log N$. The total number of samples
in an interval of length $\rho$, divided by $\log N$, differs from
$\rho$ by at most $1/\log N$.

Put $w=\min(\omega,M)$. On a full grid interval, replace both the
spatial observable and the sampled observable by their cell
representatives. The first replacement costs at most $w\rho$.
The cell-count estimate contributes at most
$M/\log N+Mm(\varphi\rho+2/\log N)$; the additional
$M/\log N$ also bounds a possible discrepancy in the total cell
mass. Replacing the representative samples back by actual samples
costs at most $w(\rho+1/\log N)$. Hence the integral is bounded above
by the sampled sum, normalized by $\log N$, plus
\[
 (2w+Mm\varphi)\rho+\frac{w+M+2Mm}{\log N}.
\]
Sum this inequality over the full grid intervals contained in
$[r,s]$. There are at most $\lfloor T/\rho\rfloor+1$ of them.
The two remaining portions have total length at most $2\rho$,
and their integrals are at most $2M\rho$. By nonnegativity, any
additional samples can be retained on the right. We obtain a
one-sided error bounded by
\[
 2wT+Mm\varphi T+2M\rho+
 \frac{(w+M+2Mm)(\lfloor T/\rho\rfloor+1)}{\log N}.
\]
Apply the same argument to $M-f$. Subtracting its integral and sum
from those of the constant $M$ gives the reverse inequality, with
an additional error at most $M/\log N$. This subtraction is valid
by the assumed temporal integrability.

Finally, the interpolated integral~\eqref{eq:empirical-integral}
differs from the sampled sum by at most $2M/\log N$: only the two
endpoint columns can be partial. Thus the absolute error is at most
the last display plus $3M/\log N$. Since $m\ge1$, $w\le M$,
$w\le\omega$, and $\lfloor T/\rho\rfloor\le T/\rho$, this is
bounded by~\eqref{eq:quadrature-bound}.
\end{proof}

\begin{lemma}[Transfer of the margins]\label{lem:margin-transfer}
Let $T>0$, and assume~\eqref{eq:scalar-integrability} for a fixed
partition. Suppose an array with $K=\lfloor T\log N\rfloor$ satisfies,
for every $0\le r\le s\le S_N$ and $2\le\theta\le4$,
\begin{equation}\label{eq:interval-transfer}
 |A^{(N)}(r,s)-A(r,s)|\le e,\qquad
 |B^{(N)}_\theta(r,s)-B_\theta(r,s)|\le e,
 \qquad e\ge0.
\end{equation}
Then
\begin{align}\label{eq:margin-transfer}
 |\beta^{(N)}(S_N)-\Lambda_0(T)|&\le e+\frac{\log k}{\log N},\nonumber\\
 |\Lambda_1^{(N)}-\Lambda_1(T)|&\le\frac{13}{2}e+
                                          \frac{14\log k}{\log N}.
\end{align}
\end{lemma}

\begin{proof}
The first estimate follows by splitting $A(0,T)$ at $S_N$ and
using~\eqref{eq:horizon-rounding} and
$0\le A(S_N,T)\le(T-S_N)\log k$. For the second, first compare
the variational expressions on the common triangle
$0\le s\le u\le S_N$. The coefficient $(\theta+3)/2$ lies in
$[3,7/2]$, so~\eqref{eq:interval-transfer} bounds their difference
by $(7/2+1+2)e=13e/2$, uniformly in $s,u,\theta$. Taking an
infimum and then a supremum preserves this absolute bound.
Lemma~\ref{lem:margin-dominance} bounds the remaining change of
horizon by $8\log k\,(T-S_N)\le8\log k/\log N$. Enlarging $8$
to $14$ gives exactly~\eqref{eq:margin-transfer}.
\end{proof}

\begin{proposition}[Uniform profile transfer]\label{prop:profile-transfer}
Fix $T,e>0$ and an arbitrary positive upper bound on $\rho$.
There exist $0<\chi<1/k$, a finite $\chi$-regular partition
$\mathcal P$, and $\rho_0>0$ below this upper bound such that, for
every $0<\rho\le\rho_0$, there is $\varphi_0>0$ with the following
property. For each $0<\varphi\le\varphi_0$ and each profile satisfying
\eqref{eq:scalar-integrability} for this same partition, all
sufficiently large $N$ and all arrays satisfying~\eqref{eq:almost-like}
obey both~\eqref{eq:interval-transfer} and~\eqref{eq:margin-transfer}.
The partition is chosen before the profile and before $N$; the
large-$N$ threshold is uniform over the cut arrays.
\end{proposition}

\begin{proof}
Choose
\[
 \chi=\min\left\{\frac1{2k},\frac e{8(T+1)}\right\}
\]
and a regular partition by Lemma~\ref{lem:regular-partition}. In
Lemma~\ref{lem:scalar-quadrature}, apply the same bound $M=4\log k$
to $a_p$ and to all $\psi_p(\theta)$, $2\le\theta\le4$, with
$\omega=\chi$. Choose $\rho_0$ no larger than the prescribed upper
bound and $e/(16(4\log k+1))$, and choose
\[
 \varphi_0=\frac e{4(4mT\log k+1)}.
\]
The first three terms on the right of~\eqref{eq:quadrature-bound}
are each at most $e/4$. After these choices, increase $N$ until its
last term is at most $e/4$. This proves the interval estimates
simultaneously for every endpoint pair and every $\theta\in[2,4]$.
Lemma~\ref{lem:margin-transfer} now applies. Every choice preceding
$N$ is independent of the cut array, as required.
\end{proof}

\begin{corollary}[Transfer for a fixed partition]\label{cor:fixed-partition-transfer}
Fix $T,e>0$. There is $\chi_0>0$ such that, for every
$0<\chi\le\chi_0$ and every fixed $\chi$-regular partition, the
conclusion of Proposition~\ref{prop:profile-transfer} holds with that
partition. More precisely, after any positive upper bound on $\rho$
is prescribed, one can choose $\rho_0$ below it, then $\varphi_0$
for each $0<\rho\le\rho_0$, and finally a common large-$N$
threshold for each $0<\varphi\le\varphi_0$. This threshold is
uniform over both profiles satisfying~\eqref{eq:scalar-integrability}
and cut arrays satisfying~\eqref{eq:almost-like}.
\end{corollary}

\begin{proof}
Take $\chi_0$ to be half the value of $\chi$ chosen in the proof of
Proposition~\ref{prop:profile-transfer}. The spatial contribution
$2T\chi$ is then at most $e/4$ for any $\chi\le\chi_0$. For the
given partition, repeat the choices of $\rho_0$, $\varphi_0$, and
$N$ in that proof. Only the number of cells enters the remaining
error terms; none involves the profile or the array.
\end{proof}

For the critical-time statement we use a consistent infinite array for
each $N$, so that finite horizons are obtained by restriction. Such
arrays admit arbitrarily fine approximation to $(\mu_t)$ if, for every
$T>0$ and every $\chi_*>0$, there are $0<\chi\le\chi_*$, a
$\chi$-regular partition satisfying~\eqref{eq:scalar-integrability},
and $g\in(0,1/k)$, $\gamma>0$, $\bar\rho>0$ such that, for every
fixed $0<\rho\le\bar\rho$ and $\varphi>0$, all sufficiently large
$N$ satisfy both~\eqref{eq:almost-like} and~\eqref{eq:local-good}
on $[0,T]$. This specifies the order of quantifiers used below.

